% Options for packages loaded elsewhere
\PassOptionsToPackage{unicode}{hyperref}
\PassOptionsToPackage{hyphens}{url}
\PassOptionsToPackage{dvipsnames,svgnames,x11names}{xcolor}
%
\documentclass[
  11pt,
  a4paper,
  DIV=11,
  numbers=noendperiod]{scrartcl}

\usepackage{amsmath,amssymb}
\usepackage{iftex}
\ifPDFTeX
  \usepackage[T1]{fontenc}
  \usepackage[utf8]{inputenc}
  \usepackage{textcomp} % provide euro and other symbols
\else % if luatex or xetex
  \usepackage{unicode-math}
  \defaultfontfeatures{Scale=MatchLowercase}
  \defaultfontfeatures[\rmfamily]{Ligatures=TeX,Scale=1}
\fi
\usepackage{lmodern}
\ifPDFTeX\else  
    % xetex/luatex font selection
    \setmainfont[]{Avenir Next}
\fi
% Use upquote if available, for straight quotes in verbatim environments
\IfFileExists{upquote.sty}{\usepackage{upquote}}{}
\IfFileExists{microtype.sty}{% use microtype if available
  \usepackage[]{microtype}
  \UseMicrotypeSet[protrusion]{basicmath} % disable protrusion for tt fonts
}{}
\makeatletter
\@ifundefined{KOMAClassName}{% if non-KOMA class
  \IfFileExists{parskip.sty}{%
    \usepackage{parskip}
  }{% else
    \setlength{\parindent}{0pt}
    \setlength{\parskip}{6pt plus 2pt minus 1pt}}
}{% if KOMA class
  \KOMAoptions{parskip=half}}
\makeatother
\usepackage{xcolor}
\setlength{\emergencystretch}{3em} % prevent overfull lines
\setcounter{secnumdepth}{-\maxdimen} % remove section numbering
% Make \paragraph and \subparagraph free-standing
\makeatletter
\ifx\paragraph\undefined\else
  \let\oldparagraph\paragraph
  \renewcommand{\paragraph}{
    \@ifstar
      \xxxParagraphStar
      \xxxParagraphNoStar
  }
  \newcommand{\xxxParagraphStar}[1]{\oldparagraph*{#1}\mbox{}}
  \newcommand{\xxxParagraphNoStar}[1]{\oldparagraph{#1}\mbox{}}
\fi
\ifx\subparagraph\undefined\else
  \let\oldsubparagraph\subparagraph
  \renewcommand{\subparagraph}{
    \@ifstar
      \xxxSubParagraphStar
      \xxxSubParagraphNoStar
  }
  \newcommand{\xxxSubParagraphStar}[1]{\oldsubparagraph*{#1}\mbox{}}
  \newcommand{\xxxSubParagraphNoStar}[1]{\oldsubparagraph{#1}\mbox{}}
\fi
\makeatother

\usepackage{color}
\usepackage{fancyvrb}
\newcommand{\VerbBar}{|}
\newcommand{\VERB}{\Verb[commandchars=\\\{\}]}
\DefineVerbatimEnvironment{Highlighting}{Verbatim}{commandchars=\\\{\}}
% Add ',fontsize=\small' for more characters per line
\usepackage{framed}
\definecolor{shadecolor}{RGB}{241,243,245}
\newenvironment{Shaded}{\begin{snugshade}}{\end{snugshade}}
\newcommand{\AlertTok}[1]{\textcolor[rgb]{0.68,0.00,0.00}{#1}}
\newcommand{\AnnotationTok}[1]{\textcolor[rgb]{0.37,0.37,0.37}{#1}}
\newcommand{\AttributeTok}[1]{\textcolor[rgb]{0.40,0.45,0.13}{#1}}
\newcommand{\BaseNTok}[1]{\textcolor[rgb]{0.68,0.00,0.00}{#1}}
\newcommand{\BuiltInTok}[1]{\textcolor[rgb]{0.00,0.23,0.31}{#1}}
\newcommand{\CharTok}[1]{\textcolor[rgb]{0.13,0.47,0.30}{#1}}
\newcommand{\CommentTok}[1]{\textcolor[rgb]{0.37,0.37,0.37}{#1}}
\newcommand{\CommentVarTok}[1]{\textcolor[rgb]{0.37,0.37,0.37}{\textit{#1}}}
\newcommand{\ConstantTok}[1]{\textcolor[rgb]{0.56,0.35,0.01}{#1}}
\newcommand{\ControlFlowTok}[1]{\textcolor[rgb]{0.00,0.23,0.31}{\textbf{#1}}}
\newcommand{\DataTypeTok}[1]{\textcolor[rgb]{0.68,0.00,0.00}{#1}}
\newcommand{\DecValTok}[1]{\textcolor[rgb]{0.68,0.00,0.00}{#1}}
\newcommand{\DocumentationTok}[1]{\textcolor[rgb]{0.37,0.37,0.37}{\textit{#1}}}
\newcommand{\ErrorTok}[1]{\textcolor[rgb]{0.68,0.00,0.00}{#1}}
\newcommand{\ExtensionTok}[1]{\textcolor[rgb]{0.00,0.23,0.31}{#1}}
\newcommand{\FloatTok}[1]{\textcolor[rgb]{0.68,0.00,0.00}{#1}}
\newcommand{\FunctionTok}[1]{\textcolor[rgb]{0.28,0.35,0.67}{#1}}
\newcommand{\ImportTok}[1]{\textcolor[rgb]{0.00,0.46,0.62}{#1}}
\newcommand{\InformationTok}[1]{\textcolor[rgb]{0.37,0.37,0.37}{#1}}
\newcommand{\KeywordTok}[1]{\textcolor[rgb]{0.00,0.23,0.31}{\textbf{#1}}}
\newcommand{\NormalTok}[1]{\textcolor[rgb]{0.00,0.23,0.31}{#1}}
\newcommand{\OperatorTok}[1]{\textcolor[rgb]{0.37,0.37,0.37}{#1}}
\newcommand{\OtherTok}[1]{\textcolor[rgb]{0.00,0.23,0.31}{#1}}
\newcommand{\PreprocessorTok}[1]{\textcolor[rgb]{0.68,0.00,0.00}{#1}}
\newcommand{\RegionMarkerTok}[1]{\textcolor[rgb]{0.00,0.23,0.31}{#1}}
\newcommand{\SpecialCharTok}[1]{\textcolor[rgb]{0.37,0.37,0.37}{#1}}
\newcommand{\SpecialStringTok}[1]{\textcolor[rgb]{0.13,0.47,0.30}{#1}}
\newcommand{\StringTok}[1]{\textcolor[rgb]{0.13,0.47,0.30}{#1}}
\newcommand{\VariableTok}[1]{\textcolor[rgb]{0.07,0.07,0.07}{#1}}
\newcommand{\VerbatimStringTok}[1]{\textcolor[rgb]{0.13,0.47,0.30}{#1}}
\newcommand{\WarningTok}[1]{\textcolor[rgb]{0.37,0.37,0.37}{\textit{#1}}}

\providecommand{\tightlist}{%
  \setlength{\itemsep}{0pt}\setlength{\parskip}{0pt}}\usepackage{longtable,booktabs,array}
\usepackage{calc} % for calculating minipage widths
% Correct order of tables after \paragraph or \subparagraph
\usepackage{etoolbox}
\makeatletter
\patchcmd\longtable{\par}{\if@noskipsec\mbox{}\fi\par}{}{}
\makeatother
% Allow footnotes in longtable head/foot
\IfFileExists{footnotehyper.sty}{\usepackage{footnotehyper}}{\usepackage{footnote}}
\makesavenoteenv{longtable}
\usepackage{graphicx}
\makeatletter
\newsavebox\pandoc@box
\newcommand*\pandocbounded[1]{% scales image to fit in text height/width
  \sbox\pandoc@box{#1}%
  \Gscale@div\@tempa{\textheight}{\dimexpr\ht\pandoc@box+\dp\pandoc@box\relax}%
  \Gscale@div\@tempb{\linewidth}{\wd\pandoc@box}%
  \ifdim\@tempb\p@<\@tempa\p@\let\@tempa\@tempb\fi% select the smaller of both
  \ifdim\@tempa\p@<\p@\scalebox{\@tempa}{\usebox\pandoc@box}%
  \else\usebox{\pandoc@box}%
  \fi%
}
% Set default figure placement to htbp
\def\fps@figure{htbp}
\makeatother

\usepackage[document]{ragged2e}
\KOMAoption{captions}{tableheading}
\makeatletter
\@ifpackageloaded{tcolorbox}{}{\usepackage[skins,breakable]{tcolorbox}}
\@ifpackageloaded{fontawesome5}{}{\usepackage{fontawesome5}}
\definecolor{quarto-callout-color}{HTML}{909090}
\definecolor{quarto-callout-note-color}{HTML}{0758E5}
\definecolor{quarto-callout-important-color}{HTML}{CC1914}
\definecolor{quarto-callout-warning-color}{HTML}{EB9113}
\definecolor{quarto-callout-tip-color}{HTML}{00A047}
\definecolor{quarto-callout-caution-color}{HTML}{FC5300}
\definecolor{quarto-callout-color-frame}{HTML}{acacac}
\definecolor{quarto-callout-note-color-frame}{HTML}{4582ec}
\definecolor{quarto-callout-important-color-frame}{HTML}{d9534f}
\definecolor{quarto-callout-warning-color-frame}{HTML}{f0ad4e}
\definecolor{quarto-callout-tip-color-frame}{HTML}{02b875}
\definecolor{quarto-callout-caution-color-frame}{HTML}{fd7e14}
\makeatother
\makeatletter
\@ifpackageloaded{caption}{}{\usepackage{caption}}
\AtBeginDocument{%
\ifdefined\contentsname
  \renewcommand*\contentsname{Table of contents}
\else
  \newcommand\contentsname{Table of contents}
\fi
\ifdefined\listfigurename
  \renewcommand*\listfigurename{List of Figures}
\else
  \newcommand\listfigurename{List of Figures}
\fi
\ifdefined\listtablename
  \renewcommand*\listtablename{List of Tables}
\else
  \newcommand\listtablename{List of Tables}
\fi
\ifdefined\figurename
  \renewcommand*\figurename{Figure}
\else
  \newcommand\figurename{Figure}
\fi
\ifdefined\tablename
  \renewcommand*\tablename{Table}
\else
  \newcommand\tablename{Table}
\fi
}
\@ifpackageloaded{float}{}{\usepackage{float}}
\floatstyle{ruled}
\@ifundefined{c@chapter}{\newfloat{codelisting}{h}{lop}}{\newfloat{codelisting}{h}{lop}[chapter]}
\floatname{codelisting}{Listing}
\newcommand*\listoflistings{\listof{codelisting}{List of Listings}}
\makeatother
\makeatletter
\makeatother
\makeatletter
\@ifpackageloaded{caption}{}{\usepackage{caption}}
\@ifpackageloaded{subcaption}{}{\usepackage{subcaption}}
\makeatother

\usepackage{bookmark}

\IfFileExists{xurl.sty}{\usepackage{xurl}}{} % add URL line breaks if available
\urlstyle{same} % disable monospaced font for URLs
\hypersetup{
  pdftitle={3. Zahlen},
  colorlinks=true,
  linkcolor={blue},
  filecolor={Maroon},
  citecolor={Blue},
  urlcolor={Blue},
  pdfcreator={LaTeX via pandoc}}


\title{3. Zahlen}
\author{}
\date{}

\begin{document}
\maketitle


\paragraph{\texorpdfstring{3.1 Ganze Zahlen -
\texttt{int}}{3.1 Ganze Zahlen - int}}\label{ganze-zahlen---int}

Die ganzen Zahlen haben in Python den Datentyp int (=interger =
ganzzahlig)

Für ganze Zahlen gibt es keine Begrenzung, sie können beliebig groß
sein.

Es gibt die bekannten Rechenoperationen wie Addition, Subtraktion,
Multiplikation und Division.

\begin{Shaded}
\begin{Highlighting}[numbers=left,,]
\DecValTok{2}
\end{Highlighting}
\end{Shaded}

\begin{verbatim}
2
\end{verbatim}

\begin{Shaded}
\begin{Highlighting}[numbers=left,,]
\DecValTok{2}\OperatorTok{+}\DecValTok{3}
\end{Highlighting}
\end{Shaded}

\begin{verbatim}
5
\end{verbatim}

\begin{Shaded}
\begin{Highlighting}[numbers=left,,]
\DecValTok{5}\OperatorTok{{-}}\DecValTok{13}
\end{Highlighting}
\end{Shaded}

\begin{verbatim}
-8
\end{verbatim}

\begin{Shaded}
\begin{Highlighting}[numbers=left,,]
\DecValTok{12}\OperatorTok{*}\DecValTok{324}
\end{Highlighting}
\end{Shaded}

\begin{verbatim}
3888
\end{verbatim}

\begin{Shaded}
\begin{Highlighting}[numbers=left,,]
\DecValTok{6}\OperatorTok{/}\DecValTok{3}
\end{Highlighting}
\end{Shaded}

\begin{verbatim}
2.0
\end{verbatim}

\paragraph{\texorpdfstring{3.2 Kommazahlen -
\texttt{float}}{3.2 Kommazahlen - float}}\label{kommazahlen---float}

Der Wert der Rechnung 6/3 ist 2.0. Es gibt folglich auch Kommawerte in
python. Diese Zahlen habe den Datentyp \texttt{float} Die Datentypen
sind nicht begrenzt. Auch die Gleitkommazahlen kennen die bekannten
Rechenoperationen wie Addition, Subtraktion, Multiplikation und
Division.

\begin{Shaded}
\begin{Highlighting}[numbers=left,,]
\DecValTok{8}\OperatorTok{/}\DecValTok{4}
\end{Highlighting}
\end{Shaded}

\begin{verbatim}
2.0
\end{verbatim}

\begin{Shaded}
\begin{Highlighting}[numbers=left,,]
\DecValTok{5}\OperatorTok{/}\DecValTok{4}
\end{Highlighting}
\end{Shaded}

\begin{verbatim}
1.25
\end{verbatim}

\begin{Shaded}
\begin{Highlighting}[numbers=left,,]
\DecValTok{7}\OperatorTok{/}\DecValTok{3}
\end{Highlighting}
\end{Shaded}

\begin{verbatim}
2.3333333333333335
\end{verbatim}

\begin{Shaded}
\begin{Highlighting}[numbers=left,,]
\FloatTok{1.2} \OperatorTok{+}\FloatTok{3.34}
\end{Highlighting}
\end{Shaded}

\begin{verbatim}
4.54
\end{verbatim}

\begin{Shaded}
\begin{Highlighting}[numbers=left,,]
\OperatorTok{{-}}\FloatTok{3.45} \OperatorTok{{-}} \FloatTok{7.4}
\end{Highlighting}
\end{Shaded}

\begin{verbatim}
-10.850000000000001
\end{verbatim}

\begin{Shaded}
\begin{Highlighting}[numbers=left,,]
\FloatTok{4.6}\OperatorTok{/}\FloatTok{4.3}
\end{Highlighting}
\end{Shaded}

\begin{verbatim}
1.069767441860465
\end{verbatim}

\begin{Shaded}
\begin{Highlighting}[numbers=left,,]
\OperatorTok{{-}}\FloatTok{2.3}\OperatorTok{*{-}}\FloatTok{3.6}
\end{Highlighting}
\end{Shaded}

\begin{verbatim}
8.28
\end{verbatim}

\begin{Shaded}
\begin{Highlighting}[numbers=left,,]
\DecValTok{1}\OperatorTok{+}\FloatTok{2.3}
\end{Highlighting}
\end{Shaded}

\begin{verbatim}
3.3
\end{verbatim}

\paragraph{\texorpdfstring{3.3 Komplexe Zahlen -
\texttt{complex}}{3.3 Komplexe Zahlen - complex}}\label{komplexe-zahlen---complex}

Komplexe Zahlen haben in Python den Datentyp \texttt{complex}. Sie
bestehen aus einem Realteil und einem Imaginärteil, die durch ein
\texttt{j} getrennt sind. Zum Beispiel:

\begin{Shaded}
\begin{Highlighting}[numbers=left,,]
\DecValTok{3}\OperatorTok{+}\OtherTok{4j}
\end{Highlighting}
\end{Shaded}

\begin{verbatim}
(3+4j)
\end{verbatim}

In Python können komplexe Zahlen direkt mit den Operatoren \texttt{+},
\texttt{-}, \texttt{*} und \texttt{/} addiert, subtrahiert,
multipliziert und dividiert werden.

\begin{itemize}
\tightlist
\item
  \textbf{Addition:} \texttt{(a+bj)\ +\ (c+dj)}
\item
  \textbf{Subtraktion:} \texttt{(a+bj)\ -\ (c+dj)}
\item
  \textbf{Multiplikation:} \texttt{(a+bj)\ *\ (c+dj)}
\item
  \textbf{Division:} \texttt{(a+bj)\ /\ (c+dj)}
\end{itemize}

Komplexe Zahlen werden in der Klassenstufe 13.2 eingeführt und im
Studium der Mathematik im Bereich der Funktionentheorie ausführlich
untersucht.

\begin{Shaded}
\begin{Highlighting}[numbers=left,,]
\NormalTok{(}\DecValTok{3}\OperatorTok{+}\OtherTok{4j}\NormalTok{) }\OperatorTok{+}\NormalTok{ (}\OperatorTok{{-}}\DecValTok{2}\OperatorTok{{-}}\OtherTok{5j}\NormalTok{)}
\end{Highlighting}
\end{Shaded}

\begin{verbatim}
(1-1j)
\end{verbatim}

\begin{Shaded}
\begin{Highlighting}[numbers=left,,]
\NormalTok{(}\DecValTok{3}\OperatorTok{+}\OtherTok{4j}\NormalTok{) }\OperatorTok{{-}}\NormalTok{ (}\OperatorTok{{-}}\DecValTok{2}\OperatorTok{{-}}\OtherTok{5j}\NormalTok{)}
\end{Highlighting}
\end{Shaded}

\begin{verbatim}
(5+9j)
\end{verbatim}

\begin{Shaded}
\begin{Highlighting}[numbers=left,,]
\NormalTok{(}\DecValTok{3}\OperatorTok{+}\OtherTok{4j}\NormalTok{) }\OperatorTok{*}\NormalTok{ (}\OperatorTok{{-}}\DecValTok{2}\OperatorTok{{-}}\OtherTok{5j}\NormalTok{)}
\end{Highlighting}
\end{Shaded}

\begin{verbatim}
(14-23j)
\end{verbatim}

\begin{Shaded}
\begin{Highlighting}[numbers=left,,]
\NormalTok{(}\DecValTok{3}\OperatorTok{+}\OtherTok{4j}\NormalTok{) }\OperatorTok{/}\NormalTok{ (}\OperatorTok{{-}}\DecValTok{2}\OperatorTok{{-}}\OtherTok{5j}\NormalTok{)}
\end{Highlighting}
\end{Shaded}

\begin{verbatim}
(-0.896551724137931+0.24137931034482757j)
\end{verbatim}

\paragraph{\texorpdfstring{3.4 Logische Werte -
\texttt{bool}}{3.4 Logische Werte - bool}}\label{logische-werte---bool}

Computer kennt zwei Zustände: Strom an und Strom aus. Diese ist also die
kleinste Informationseinheit, die in einem Computer gespeichert werden
kann und wird als binary digit (binär Ziffer), oder kurz Bit,
bezeichnet. In Python gibt es den Datentyp \texttt{bool}, der zwei
mögliche Werte hat: \texttt{True} (wahr) und \texttt{False} (falsch).

\begin{Shaded}
\begin{Highlighting}[numbers=left,,]
\VariableTok{True}

\end{Highlighting}
\end{Shaded}

\begin{verbatim}
True
\end{verbatim}

\begin{Shaded}
\begin{Highlighting}[numbers=left,,]
\VariableTok{False}
\end{Highlighting}
\end{Shaded}

\begin{verbatim}
False
\end{verbatim}

\begin{Shaded}
\begin{Highlighting}[numbers=left,,]
\DecValTok{5}\OperatorTok{\textgreater{}}\DecValTok{10}
\end{Highlighting}
\end{Shaded}

\begin{verbatim}
False
\end{verbatim}

\begin{Shaded}
\begin{Highlighting}[numbers=left,,]
\DecValTok{5}\OperatorTok{==}\DecValTok{5}
\end{Highlighting}
\end{Shaded}

\begin{verbatim}
True
\end{verbatim}

\paragraph{3.5 Grundoperationen auf
Bits}\label{grundoperationen-auf-bits}

Mittels der logischen Interpretation von Bits können drei
Grundoperationen definiert werden, die auf Bits angewendet werden
können. Diese Operationen sind die Basis für alle weiteren logischen
Operationen und ermöglichen die Manipulation von Bits in der Informatik.

\textbf{1. Logisches AND: \(b_1 \land b_2\)}

\begin{longtable}[]{@{}ccc@{}}
\toprule\noalign{}
\(b_1\) & \(b_2\) & \(b_1 \land b_2\) \\
\midrule\noalign{}
\endhead
\bottomrule\noalign{}
\endlastfoot
0 & 0 & 0 \\
1 & 0 & 0 \\
0 & 1 & 0 \\
1 & 1 & 1 \\
\end{longtable}

\textbf{2. Logisches OR: \(b_1 \lor b_2\)}

\begin{longtable}[]{@{}ccc@{}}
\toprule\noalign{}
\(b_1\) & \(b_2\) & \(b_1 \lor b_2\) \\
\midrule\noalign{}
\endhead
\bottomrule\noalign{}
\endlastfoot
0 & 0 & 0 \\
1 & 0 & 1 \\
0 & 1 & 1 \\
1 & 1 & 1 \\
\end{longtable}

\textbf{3. Logisches NOT (Negation, Komplement): \(\neg b\)}

\begin{longtable}[]{@{}cc@{}}
\toprule\noalign{}
\(b\) & \(\neg b\) \\
\midrule\noalign{}
\endhead
\bottomrule\noalign{}
\endlastfoot
0 & 1 \\
1 & 0 \\
\end{longtable}

\begin{tcolorbox}[enhanced jigsaw, colbacktitle=quarto-callout-important-color!10!white, leftrule=.75mm, colback=white, toprule=.15mm, arc=.35mm, titlerule=0mm, opacityback=0, breakable, coltitle=black, bottomtitle=1mm, opacitybacktitle=0.6, rightrule=.15mm, title=\textcolor{quarto-callout-important-color}{\faExclamation}\hspace{0.5em}{Satz}, toptitle=1mm, bottomrule=.15mm, left=2mm, colframe=quarto-callout-important-color-frame]

Mit den drei Grundoperationen \textbf{AND}, \textbf{OR} und \textbf{NOT}
können alle möglichen Operationen auf Bits definiert werden.

\end{tcolorbox}

\textbf{Beispiel 1}

\begin{longtable}[]{@{}ccc@{}}
\toprule\noalign{}
\(b_1\) & \(b_2\) & \(f(b_1, b_2)\) \\
\midrule\noalign{}
\endhead
\bottomrule\noalign{}
\endlastfoot
0 & 0 & 0 \\
1 & 0 & 0 \\
0 & 1 & 0 \\
1 & 1 & 1 \\
\end{longtable}

Auflösung: \(f(b_1,b_2)=b_1 \land b_2\)

\textbf{Beispiel 2}

\begin{longtable}[]{@{}ccc@{}}
\toprule\noalign{}
\(b_1\) & \(b_2\) & \(f(b_1, b_2)\) \\
\midrule\noalign{}
\endhead
\bottomrule\noalign{}
\endlastfoot
0 & 0 & 1 \\
1 & 0 & 1 \\
0 & 1 & 0 \\
1 & 1 & 1 \\
\end{longtable}

\textbf{Bemerkung:}

entweder oder

\begin{longtable}[]{@{}ccc@{}}
\toprule\noalign{}
\(b_1\) & \(b_2\) & \(b_1 \text{ xor }b_2\) \\
\midrule\noalign{}
\endhead
\bottomrule\noalign{}
\endlastfoot
0 & 0 & 0 \\
1 & 0 & 1 \\
0 & 1 & 1 \\
1 & 1 & 0 \\
\end{longtable}

\subsubsection{Exkurs Aussagenlogik}\label{exkurs-aussagenlogik}

Die Aussagenlogik ist ein grundlegendes Werkzeug, um Sachverhalte
präzise zu beschreiben und logisch zu analysieren. In der Mathematik
wird sie genutzt, um Beweise zu führen, Aussagen zu verknüpfen und die
Gültigkeit von Argumenten zu überprüfen. In der Informatik bildet die
Aussagenlogik die Basis für die Entwicklung von Algorithmen, die
Funktionsweise von Schaltungen, die Programmierung von Bedingungen (if,
while, etc.) und die Fehleranalyse. Sie hilft, komplexe Probleme in
einfache, logisch überprüfbare Einzelschritte zu zerlegen.

\paragraph{Grundbegriffe}\label{grundbegriffe}

\begin{itemize}
\tightlist
\item
  \textbf{Aussage:} Ein Satz, der eindeutig wahr oder falsch ist.
\item
  \textbf{Wahrheitswert:} Jede Aussage ist entweder wahr (W) oder falsch
  (F).
\item
  \textbf{Logische Operatoren:}

  \begin{itemize}
  \tightlist
  \item
    \textbf{Negation (¬):} Nicht
  \item
    \textbf{Konjunktion (∧):} Und
  \item
    \textbf{Disjunktion (∨):} Oder
  \item
    \textbf{Implikation (→):} Wenn \ldots{} dann
  \item
    \textbf{Äquivalenz (↔):} Genau dann, wenn
  \end{itemize}
\end{itemize}

\paragraph{Wahrheitstafeln}\label{wahrheitstafeln}

\begin{longtable}[]{@{}lllllll@{}}
\toprule\noalign{}
A & B & ¬A & A ∧ B & A ∨ B & A → B & A ↔ B \\
\midrule\noalign{}
\endhead
\bottomrule\noalign{}
\endlastfoot
W & W & F & W & W & W & W \\
W & F & F & F & W & F & F \\
F & W & W & F & W & W & F \\
F & F & W & F & F & W & W \\
\end{longtable}

\paragraph{Beispiele}\label{beispiele}

\begin{itemize}
\tightlist
\item
  \textbf{Negation:} Die Aussage „Es regnet`` wird zu „Es regnet
  nicht``.
\item
  \textbf{Konjunktion:} „Es regnet und es ist warm.``
\item
  \textbf{Disjunktion:} „Es regnet oder es ist warm.``
\item
  \textbf{Implikation:} „Wenn es regnet, dann ist die Straße nass.``
\item
  \textbf{Äquivalenz:} „Die Lampe leuchtet genau dann, wenn der Schalter
  an ist.``
\end{itemize}

\paragraph{Aufgaben:}\label{aufgaben}

\begin{enumerate}
\def\labelenumi{\arabic{enumi}.}
\tightlist
\item
  Erstelle die Wahrheitstafel für die Aussage: ¬(A ∨ B)
\item
  Überprüfe, ob die Aussage (A ∧ B) → A immer wahr ist.
\item
  Formuliere eigene Beispiele für die logischen Operatoren.
\end{enumerate}




\end{document}
